\section{2013级工科数学分析下 A卷}

\subsection{资料来源}
\begin{itemize}
  \item 试题：2013 级软件工科数学分析下 A 卷，已导出页面图校对。
  \item 答案：同卷附解答文件。
\end{itemize}

\subsection{试题}

\paragraph*{试卷信息}
诚信应考，考试作弊将带来严重后果！

华南理工大学本科生期末考试《工科数学分析》2013--2014学年第二学期期末考试试卷（A）卷。

注意事项：开考前请将密封线内各项信息填写清楚；所有答案请直接答在试卷上（或答题纸上）；考试形式：闭卷；本试卷共5个大题，满分100分，考试时间120分钟。

\paragraph*{一、单项选择题（每小题3分，共15分）}
\begin{enumerate}
  \item 若 $y_1,y_2$ 是非齐次线性方程 $y''+ay'+by=f(x)$ 的两个特解，则下面结论中正确的是（\quad）。
  \begin{enumerate}
    \item[A.] $y_1+y_2$ 是非齐次线性方程的解。
    \item[B.] $y_1-y_2$ 是非齐次线性方程的解。
    \item[C.] $y_1+y_2$ 是 $y''+ay'+by=0$ 的解。
    \item[D.] $y_1-y_2$ 是 $y''+ay'+by=0$ 的解。
  \end{enumerate}

  \item 设 $z=f(x,y)$ 可微，且 $f(x,x^2)=\frac12,\ \left.f_x(x,y)\right|_{y=x^2}=\frac1x$，则 $\left.f_y(x,y)\right|_{y=x^2}$（\quad）。
  \begin{enumerate}
    \item[A.] $-\dfrac{1}{x^2}$
    \item[B.] $\dfrac{1}{x^2}$
    \item[C.] $\dfrac{1}{2x^2}$
    \item[D.] $-\dfrac{1}{2x^2}$
  \end{enumerate}

  \item 设 $C:x^2+y^2=a^2$ 取逆时针方向，则 $\displaystyle\oint_C \frac{(x+y)\,dx-(x-y)\,dy}{x^2+y^2}$（\quad）。
  \begin{enumerate}
    \item[A.] $-2\pi$
    \item[B.] $2\pi$
    \item[C.] $0$
    \item[D.] $-\pi$
  \end{enumerate}

  \item 幂级数 $\displaystyle\sum_{n=1}^{\infty}\frac{(-1)^{n-1}}{n}(x-1)^n$ 的收敛域是（\quad）。
  \begin{enumerate}
    \item[A.] $(0,2)$
    \item[B.] $(0,2]$
    \item[C.] $[0,2)$
    \item[D.] $[0,2]$
  \end{enumerate}

  \item 设 $f(x)=x^2\ (0\le x<1)$，而
  \[
    S(x)=\sum_{n=1}^{\infty}b_n\sin n\pi x\quad(-\infty<x<+\infty),
    \qquad
    b_n=2\int_0^1 f(x)\sin n\pi x\,dx.
  \]
  则 $S\!\left(-\dfrac12\right)$（\quad）。
  \begin{enumerate}
    \item[A.] $-\dfrac12$
    \item[B.] $\dfrac14$
    \item[C.] $-\dfrac14$
    \item[D.] $\dfrac12$
  \end{enumerate}
\end{enumerate}

\paragraph*{二、填空题（每小题3分，共15分）}
\begin{enumerate}
  \item 微分方程 $y^{(4)}-y=0$ 的通解为 \underline{\hspace{5cm}}。
  \item 函数 $z=2x^2+y^2$ 在点 $(1,1)$ 处沿该点梯度方向的方向导数为 \underline{\hspace{5cm}}。
  \item 交换积分的顺序，则 $\displaystyle\int_1^e dx\int_0^{\ln x} f(x,y)\,dy=$ \underline{\hspace{5cm}}。
  \item 设 $\Sigma:x^2+y^2+z^2=a^2$，则 $\displaystyle\iint_{\Sigma}(x^2+y^2+z^2)\,dS=$ \underline{\hspace{5cm}}。
  \item 设幂级数 $\displaystyle\sum_{n=1}^{\infty}a_nx^n$ 的收敛半径为 $3$，则幂级数
  \[
    \sum_{n=1}^{\infty}n a_n(x-1)^{n+1}
  \]
  的收敛区间为 \underline{\hspace{5cm}}。
\end{enumerate}

\paragraph*{三、计算题（每小题10分，共50分）}
\begin{enumerate}
  \item 设函数 $f(u)$ 具有二阶连续导数，而 $z=f(e^x\sin y)$ 满足方程 $\dfrac{\partial^2 z}{\partial x^2}+\dfrac{\partial^2 z}{\partial y^2}=e^{2x}z$，求 $f(u)$。
  \item 计算三重积分 $I=\displaystyle\iiint_{\Omega}(x+z)\,dV$，其中 $\Omega$ 是由曲面 $z=x^2+y^2$ 与曲面 $z=1-x^2-y^2$ 所围成的区域。
  \item 设 $f(x)$ 在 $(-\infty,+\infty)$ 内有连续导数，$L$ 是从点 $A(3,\frac23)$ 到点 $B(1,2)$ 的直线段，计算曲线积分
  \[
    I=\int_L \frac{1+y^2f(xy)}{y}\,dx
    +\frac{x}{y^2}\bigl[y^2f(xy)-1\bigr]\,dy.
  \]
  \item 设 $\Sigma$ 是锥面 $z=\sqrt{x^2+y^2}$ 被平面 $z=0$ 及 $z=1$ 所截部分的下侧，计算第二类曲面积分 $I=\displaystyle\iint_{\Sigma}x\,dy\,dz+y\,dz\,dx+(z^2-2z)\,dx\,dy$。
  \item 求幂级数 $\displaystyle\sum_{n=1}^{\infty}\frac{2n+1}{n!}x^{2n}$ 的和函数。
\end{enumerate}

\paragraph*{四、证明题（本题10分）}
\begin{enumerate}
  \setcounter{enumi}{5}
  \item 证明级数 $\displaystyle\sum_{n=1}^{\infty}\frac{(-1)^n(1-e^{-nx})}{n^2+x^2}$ 在 $[0,+\infty)$ 上一致收敛。
\end{enumerate}

\paragraph*{五、应用题（本题10分）}
\begin{enumerate}
  \setcounter{enumi}{6}
  \item 在第一象限内作椭球面 $\dfrac{x^2}{a^2}+\dfrac{y^2}{b^2}+\dfrac{z^2}{c^2}=1$ 的切平面，使切平面与三个坐标面所围成的四面体的体积最小，求切点的坐标。
\end{enumerate}

\subsection{答案与解析}

\paragraph*{一、单项选择题}
\begin{enumerate}
  \item 两个非齐次方程特解之差满足对应齐次方程，即 $y_1-y_2$ 是 $y''+ay'+by=0$ 的解。选 D。
  \item 对恒等式 $f(x,x^2)=\frac12$ 求导，得 $f_x(x,x^2)+2xf_y(x,x^2)=0$，代入 $f_x(x,x^2)=\frac1x$，得 $f_y(x,x^2)=-\dfrac1{2x^2}$。选 D。
  \item 在圆周上分解
  \(\displaystyle \frac{(x+y)\,dx-(x-y)\,dy}{x^2+y^2}
  =d\ln r-d\theta\)，闭曲线第一项积分为 $0$，逆时针一周第二项为 $-2\pi$。选 A。
  \item 令 $t=x-1$，该幂级数为 \(\displaystyle \ln(1+t)\)，其收敛区间为 $-1<t\le1$，即 $0<x\le2$。选 B。
  \item 该正弦级数对应 $[0,1]$ 上 $f(x)=x^2$ 的奇延拓，故 $S(-1/2)=-f(1/2)=-\dfrac14$。选 C。
\end{enumerate}

\paragraph*{二、填空题}
\begin{enumerate}
  \item 微分方程 $y^{(4)}-y=0$ 的通解为 $y=C_1e^x+C_2e^{-x}+C_3\cos x+C_4\sin x$。
  \item 函数 $z=2x^2+y^2$ 在点 $(1,1)$ 处沿该点梯度方向的方向导数为 $2\sqrt5$。
  \item 交换积分的顺序，则 $\displaystyle\int_1^e dx\int_0^{\ln x} f(x,y)\,dy$ $=\displaystyle\int_0^1dy\int_{e^y}^{e}f(x,y)\,dx$。
  \item 设 $\Sigma:x^2+y^2+z^2=a^2$，则 $\displaystyle\iint_{\Sigma}(x^2+y^2+z^2)\,dS$ $=4\pi a^4$。
  \item 设幂级数 $\displaystyle\sum_{n=1}^{\infty}a_nx^n$ 的收敛半径为 $3$，则幂级数 $\displaystyle\sum_{n=1}^{\infty}n a_n(x-1)^{n+1}$ 的收敛区间为 $(-2,4)$。
\end{enumerate}

\paragraph*{三、计算题}
\begin{enumerate}
  \item 解：由 $u=e^x\sin y$；又由 $z=f(u)$ 满足方程
  \[
    \frac{\partial^2 z}{\partial x^2}
    +\frac{\partial^2 z}{\partial y^2}=e^{2x}z,
  \]
  得到 $f''(u)=f(u)$，解得 $f(u)=C_1e^u+C_2e^{-u}$。
  \item 解：根据对称性，$\displaystyle\iiint_{\Omega}x\,dV=0$（单独2分）。利用柱面坐标可得
  \[
    I=2\pi\int_0^{1/\sqrt2}dr\int_{r^2}^{1-r^2}zr\,dz
    =\frac{\pi}{8};
  \]
  或利用球面坐标可得 $I=\dfrac{\pi}{8}$。
  \item 解：
  \[
    d\left(F(xy)+\frac{x}{y}\right)
    =\frac{1+y^2f(xy)}{y}\,dx
    +\frac{x}{y^2}\bigl[y^2f(xy)-1\bigr]\,dy.
  \]
  法一：由原积分为全微分积分（6分），则
  \[
    I=\left[F(xy)+\frac{x}{y}\right]_{(3,2/3)}^{(1,2)}=-4.
  \]
  法二：记 $F'(t)=f(t)$，则 $I=-4$。
  \item 解：法一：取参数
  \[
    x=z\cos\theta,\qquad y=z\sin\theta,\qquad
    0\le z\le1,\quad 0\le\theta\le2\pi.
  \]
  下侧对应面积向量 $(z\cos\theta,z\sin\theta,-z)\,d\theta dz$，故
  \[
    I=\int_0^{2\pi}\int_0^1(3z^2-z^3)\,dz\,d\theta
    =\frac{3\pi}{2}.
  \]
  法二：利用参数曲面积分同样可得 $I=\dfrac{3\pi}{2}$。法三：补上下底面后利用 Gauss 公式可得同一结果。
  \item 解：
  \[
    \sum_{n=1}^{\infty}\frac{2n+1}{n!}x^{2n}
    =2\sum_{n=1}^{\infty}\frac{n(x^2)^n}{n!}
    +\sum_{n=1}^{\infty}\frac{(x^2)^n}{n!}
    =2x^2e^{x^2}+e^{x^2}-1.
  \]
  因此 $S(x)=(2x^2+1)e^{x^2}-1$。
\end{enumerate}

\paragraph*{四、证明题}
证明：因为
\[
  \left|\frac{(-1)^n(1-e^{-nx})}{n^2+x^2}\right|\le \frac{1}{n^2},
\]
且 $\displaystyle\sum_{n=1}^{\infty}\frac1{n^2}$ 收敛，所以级数
\[
  \sum_{n=1}^{\infty}\frac{(-1)^n(1-e^{-nx})}{n^2+x^2}
\]
在 $[0,+\infty)$ 上一致收敛（10分）。

\paragraph*{五、应用题}
解：设 $P(x_0,y_0,z_0)$ 为椭球面上在第一象限的一点。过此点的切平面方程为
\[
  \frac{x_0x}{a^2}+\frac{y_0y}{b^2}+\frac{z_0z}{c^2}=1,
\]
即
\[
  \frac{x}{a^2/x_0}+\frac{y}{b^2/y_0}+\frac{z}{c^2/z_0}=1.
\]
此切平面与坐标面围成四面体的体积为
\[
  V=\frac{a^2b^2c^2}{6x_0y_0z_0}.
\]
（3分）
要求 $V$ 在条件
\[
  \frac{x_0^2}{a^2}+\frac{y_0^2}{b^2}+\frac{z_0^2}{c^2}=1
\]
下的最小值，只需求 $x_0y_0z_0$ 在同一条件下的最大值。由拉格朗日乘数法，只需求函数
\[
  F=x_0y_0z_0+\lambda\left(\frac{x_0^2}{a^2}
  +\frac{y_0^2}{b^2}+\frac{z_0^2}{c^2}-1\right)
\]
的驻点（6分），即 $F_{x_0}=F_{y_0}=F_{z_0}=0$（8分）。由驻点方程两两相除得
\[
  \frac{x_0}{a}=\frac{y_0}{b}=\frac{z_0}{c}.
\]
由此得
\[
  \frac{x_0^2}{a^2}=\frac{y_0^2}{b^2}=\frac{z_0^2}{c^2}=\frac13,
\]
所以
\[
  x_0=\frac{a}{\sqrt3},\qquad
  y_0=\frac{b}{\sqrt3},\qquad
  z_0=\frac{c}{\sqrt3}.
\]
（9分）
由驻点唯一，又根据实际意义最小体积存在，所以切点坐标为
\[
  \left(\frac{a}{\sqrt3},\frac{b}{\sqrt3},\frac{c}{\sqrt3}\right).
\]
（10分）
